% from ega4-2-fr.tex
\oldpage[IV-2]{65}

désignant le sous-préschéma réduit de $X$ ayant $X_0$ pour espace sous-jacent
(I, 5.2.2)~; il suffit d'appliquer (4.5.13) à $g$.

Pour (ii), en considérant un préschéma ayant $Y$ pour espace sous-jacent, et
remarquant que $Y$ est \emph{a fortiori} géométriquement connexe, il suffit
d'appliquer (i) à l'injection canonique $Y\to X$.
\end{proof}

\begin{corollary}[4.5.14]
\label{IV.2.4.5.14-fr}
Soit $X$ un $k$-préschéma. Si $x$ est un point de $X$ tel que $k(x)$ soit une
extension primaire de $k$ (en particulier si $x$ est un point rationnel de
$X$), la composante connexe de $x$ dans $X$ est géométriquement connexe.
\end{corollary}

\begin{proof}
Il suffit d'appliquer (4.5.13.1) à $Y=\operatorname{Spec}(k(x))$ et au
morphisme canonique $Y\to X$, en tenant compte de (4.5.9), qui entraîne que
$Y$ est géométriquement irréductible.
\end{proof}

\begin{proposition}[4.5.15]
\label{IV.2.4.5.15-fr}
Soient $X$ un $k$-préschéma, $x$ un point de $X$, $k'$ la fermeture algébrique
séparable de $k$ dans $k(x)$. On suppose vérifiées les conditions suivantes~:
\begin{enumerate}
\item[(i)] $X$ est connexe.
\item[(ii)] $k'$ est une extension finie de $k$.
\end{enumerate}
Alors le nombre géométrique des composantes connexes de $X$ est
$\leqslant[k':k]$, et si $k''$ est une extension galoisienne finie de $k$
contenant $k'$, les composantes connexes de $X\otimes_k k''$ sont
géométriquement connexes.
\end{proposition}

On notera que la condition (ii) est vérifiée si $X$ est \emph{localement de
type fini} sur $k$.

\begin{proof}
Il existe des extensions galoisiennes finies $k''$ de $k$ contenant $k'$ en
vertu de la condition (ii). Posons $X''=X\otimes_k k''$~; le morphisme
projection $p:X''\to X$ est fidèlement plat et fini, donc fermé (II, 6.1.10) et
par suite (2.3.6, (ii)) l'image par $p$ de toute composante connexe $X''_\alpha$
de $X''$ est égale à $X$~; $X''_\alpha$ contient donc un point
$x''_\alpha\in p^{-1}(x)$. Mais $p^{-1}(x)$ a un nombre de points égal au
nombre des composantes irréductibles de
$\operatorname{Spec}(k(x)\otimes_k k'')$ (I, 3.4.9), c'est-à-dire à $[k':k]$
(4.3.4), et \emph{a fortiori} le nombre des composantes connexes de $X''$ est
$\leqslant[k':k]$. Pour tout $x''_\alpha\in p^{-1}(x)$, $k(x''_\alpha)$ est une
extension primaire de $k''$ en vertu de (4.3.5) et de (I, 3.4.9). On peut donc
appliquer à $x''_\alpha$ et $X''$ le corollaire (4.5.14), qui prouve que toutes
les composantes connexes de $X''$ sont géométriquement connexes.
\end{proof}

\begin{corollary}[4.5.16]
\label{IV.2.4.5.16-fr}
Soient $X$ un $k$-préschéma, $(X_\alpha)$ la famille de ses composantes
connexes, et pour tout $\alpha$, soit $x_\alpha\in X_\alpha$. On suppose
vérifiées les conditions suivantes~:
\begin{enumerate}
\item[(i)] La famille $(X_\alpha)$ est finie.
\item[(ii)] Pour tout $\alpha$, la fermeture algébrique séparable $k'_\alpha$
de $k$ dans $k(x_\alpha)$ est une extension finie de $k$.
\end{enumerate}
Alors le nombre géométrique de composantes connexes de $X$ est au plus
$\sum_\alpha[k'_\alpha:k]$, et il existe une extension finie séparable $k''$ de
$k$ telle que toutes les composantes connexes de $X\otimes_k k''$ soient
géométriquement connexes.
\end{corollary}

\begin{proof}
Notant que les composantes connexes de $X$ sont ouvertes en vertu de la
condition (i), il suffit d'appliquer à chacun des sous-préschémas de $X$ induits
sur les ouverts $X_\alpha$ le résultat de (4.5.15)~; pour avoir une extension
$k''$ répondant à la question, il suffit de prendre une extension galoisienne
finie de $k$ contenant tous les $k'_\alpha$.
\end{proof}

